\section{How sure are we of each claim?}
\label{sec:confidence}

\paragraph{What was varied, and how much of it.}

\begin{center}
\footnotesize
\setlength{\tabcolsep}{4pt}
\begin{tabular}{@{}p{1.7cm}p{3.6cm}p{3.5cm}p{2.2cm}p{2.3cm}@{}}
\toprule
Experiment & Random inputs varied & Replicates & Fixed & Metric \\
\midrule
Predictive
 & 4 function draws $\times$ 8 designs
 & 32 fits per map, per $n$
 & $d=6$, $\ell=0.5$
 & CRPS, RMSE \\
 & $n \in \{5,10,20,40\}$ training pts
 & $\times$ 4 sizes $\times$ 9 maps $=1152$
 & 512 test pts
 & \\
\addlinespace
Regret
 & 4 function draws $\times$ 10 designs
 & 40 BO runs per arm per map
 & budget 40, $n_0=5$
 & final $\log_{10}$ \\
 & (confounded, see \S4)
 & $\times$ 2 arms $\times$ 9 maps $=720$
 & shared designs
 & regret \\
\addlinespace
Injectivity $\bar n$
 & ---
 & 4000 probe points, 20001-pt grid
 & draw 0
 & count of preimages \\
\addlinespace
Exact posteriors
 & --- (deterministic)
 & no sampling anywhere
 & $n=8$--12
 & Kolmogorov dist.\ \\
\bottomrule
\end{tabular}
\end{center}

\noindent A ``function draw'' is a fresh random inner function $h$; a
``design'' is a fresh set of training inputs. Varying both matters: with one
draw, a conclusion could be a property of that single sampled $h$ and the
error bars would never reveal it.

\paragraph{Each claim, with its interval.} Intervals are 95\%; paired rows use
a $t$ interval on the per-fit differences, so the pairing across the shared
design is exploited.

\begin{center}
\footnotesize
\begin{tabular}{p{6.3cm}lll}
\toprule
Claim & Estimate & 95\% CI & Confidence \\
\midrule
\multicolumn{4}{l}{\emph{Proved, then confirmed numerically --- no sampling involved}}\\
Bijective $g$: Models 2 and 3 coincide
 & $0.000000$ & exact & \textbf{certain} \\
Linear $g$: Models 1 and 3 coincide
 & $0.000000$ & exact & \textbf{certain} \\
\addlinespace
\multicolumn{4}{l}{\emph{Strong: large $t$, every replicate agrees}}\\
Composite predicts better than Direct on \emph{every} map
 & --- & --- & \textbf{strong} (9/9 maps, $t\ge3.6$) \\
\quad e.g.\ \texttt{root2}, 32 fits, 32/32 wins
 & $+0.0263$ & $[+0.0229,+0.0298]$ & $t=15.6$ \\
Bijective maps still gain (so the gain is not only information)
 & $+0.0189$ & $[+0.0143,+0.0234]$ & $t=8.5$ (\texttt{mono3}) \\
The linear map gains least, at every $n$
 & $+0.0050$ & $[+0.0035,+0.0066]$ & \textbf{strong} (4/4 sizes) \\
Gain does not shrink as $n$ grows $5\to40$
 & --- & --- & \textbf{strong} (all 9 maps) \\
\addlinespace
\multicolumn{4}{l}{\emph{Unresolved: the interval is too wide to call}}\\
Degree-3 matched pair: effect of bijectivity
 & $+0.0033$ & $[-0.0057,+0.0123]$ & $t=0.72$ \ \textbf{null} \\
Degree-5 matched pair: effect of bijectivity
 & $+0.0767$ & $[-0.0584,+0.2118]$ & $t=1.11$ \ \textbf{null} \\
Across maps, gain correlates with $\bar n$
 & $r=+0.78$ & $[+0.25,+0.95]$ & \textbf{fragile} (see below) \\
\texttt{root5}'s gain
 & $+0.1133$ & $[-0.0258,+0.2525]$ & $t=1.7$ \ \textbf{unreliable} \\
\addlinespace
\multicolumn{4}{l}{\emph{Confounded by construction --- reported, not relied on}}\\
Regret gain, \texttt{root1} (linear), 40 runs
 & $+0.260$ & $[-0.107,+0.628]$ & n.s. \\
Regret gain, \texttt{root3}, 40 runs
 & $+0.808$ & $[+0.411,+1.205]$ & see \S4 \\
\bottomrule
\end{tabular}
\end{center}

\paragraph{Why the correlation is called fragile.} It rests on one point.

\begin{center}
\footnotesize
\begin{tabular}{llll}
\toprule
Subset & Pearson $r$ & 95\% CI & Spearman \\
\midrule
All 9 maps & $+0.783$ & $[+0.247,+0.952]$ & $+0.700$ \\
Dropping \texttt{root5} (whose own CI spans zero) & $+0.558$ & $[-0.242,+0.906]$ & $+0.571$ \\
Polynomials only & $+0.766$ & $[+0.030,+0.963]$ & $+0.857$ \\
\bottomrule
\end{tabular}
\end{center}

\noindent Removing a single map moves the interval across zero. And $\bar n$
is itself correlated with degree ($r=+0.48$ across the seven polynomials), so
with nine maps the two cannot be separated. Most tellingly: \textbf{among the
five maps with $\bar n=1.00$ exactly, the gain still spans a factor of
$7.3$} ($0.0050$ to $0.0367$). Injectivity held perfectly fixed, something
else moves the gain sevenfold.

\paragraph{What the null results could and could not have seen.} The degree-3
interval has half-width $0.0090$ against a gain of about $0.020$, so it would
only have detected a bijectivity effect larger than roughly \textbf{45\% of
the gain itself}. The observed $17\%$ difference is well inside the noise.
Resolving it needs the standard error down by $2.7\times$, which is
$\mathbf{\sim\!240}$ \textbf{fits per map} instead of 32 --- about eight times
the compute, and still only two matched pairs.

\paragraph{The honest summary.} We can say with confidence that the gain is
real everywhere, that it is smallest for the linear map, that it survives on
bijective maps, and that it does not wash out with data. We \emph{cannot}
currently say whether bijectivity has a modest secondary effect: the
controlled comparison is null but underpowered, and the across-map correlation
is positive but fragile and confounded with degree.
